%%%%%%%%%%%%%%%%%%%%%%%%%%%%%%%%%%%%%%%%%%%%%%%%%%%%%%%
% % Lines starting with % are comments, which are ignored.
% % This is a handy way of indicating the date and version of
% % your document, to wit:
% %
% % LaTeX sample file
% % Modified March, 2002
% %
%%%%%%%%%%%%%%%%%%%%%%%%%%%%%%%%%%%%%%%%%%%%%%%%%%%%%%%
\DocumentMetadata{
  % PDF standards: ONE key whose value is a comma-separated list --
  % UA-2 (accessibility) plus A-4f (archival; the archival companion of
  % UA-2). Do NOT use separate pdfstandard=... keys: with ordinary
  % key-value parsing only the last one takes effect.
  pdfstandard={UA-2,A-4f},
  lang=en-US,
  tagging=on,
  tagging-setup={
    math/setup=mathml-SE
  }
}

\documentclass{ltx-talk}
% lecture-macros.tex -- shared package setup for lecture-template.tex
% (ltx-talk slides / handouts).
%
% NOTE ON PATHS: every template in this repository is stored one level up
% from where it is compiled (see README.md). All relative \input paths
% below are resolved against the folder the document is compiled in,
% which sits one directory above this file.

% --- Tables, bookmarks, strikethrough ------------------------------------
% booktabs: professional table rules (\toprule, \midrule, \bottomrule)
% bookmark: PDF outline entries (PDF bookmarks)
% ulem[normalem]: \sout strikethrough that survives math mode
\usepackage{booktabs}
\usepackage{bookmark}
\usepackage[normalem]{ulem}
\usepackage{multicol}

% --- Math ------------------------------------------------------------
% amsmath: equation environments and numbering
% mathrsfs: \mathfrak  |  stmaryrd: extra symbols
\usepackage{amsmath}
\usepackage{mathrsfs}
\usepackage{stmaryrd}

% --- Layout -------------------------------------------------------------
% float: [H] float placement
% fontspec: Unicode/OpenType fonts (compile with XeLaTeX or LuaLaTeX)
\usepackage{float}
\usepackage{fontspec}

% --- Colors ---------------------------------------------------------------
% dvipsnames: extended named color palette (the Accessible* colors below
% are built on it);
% table: \rowcolors / \cellcolor for shading tables.
% NOTE: xcolor may only be loaded ONCE per document -- a second
% \usepackage{xcolor} with different options is a fatal "Option clash"
% error, so ALL of its options are listed here.
\usepackage[dvipsnames,table]{xcolor}

% --- Multimedia -----------------------------------------------------------
% animate: embeds audio/video (and legacy Flash) media in a frame.
\usepackage{animate}

% --- Bibliography ----------------------------------------------------------
% biblatex (biber backend): citations for lecture notes.
\usepackage[
  backend=biber,
  style=alphabetic,
  url=true,
  doi=true,
  isbn=true,
  backref=true
]{biblatex}

% --- Colored boxes ----------------------------------------------------------
% talkthemetcolorbox: the colored-box theme for ltx-talk slides
% (shadow/rounded/shaded blocks, showtitle, titlepage styles).
\usepackage[
  shadow,
  rounded,
  shaded,
  blocks,
  showtitle,
  titlepage
]{talkthemetcolorbox}

% --- "theindex" helper environment ------------------------------------------
% Provides a lightweight theindex environment (an itemize) and the
% \@idxitem/\subitem/\subsubitem/\indexspace macros that LaTeX's standard
% theindex machinery expects, so hand-written index pages indent correctly.
\makeatletter
\newenvironment{theindex}
  {\begin{itemize}}
  {\end{itemize}}
\providecommand\@idxitem{\par\hangindent 40\p@}
\providecommand\subitem{\@idxitem \hspace*{20\p@}}
\providecommand\subsubitem{\@idxitem \hspace*{30\p@}}
\providecommand\indexspace{\par \vskip 10\p@ \@plus5\p@ \@minus3\p@\relax}
\makeatother

% --- Index -----------------------------------------------------------------
% imakeidx: builds the printed index.
% The redefinition below appends "|hyperpage" to every index entry so that,
% in the PDF, each index entry is a link to its page.
\usepackage{imakeidx}
\makeatletter
\let\imki@orig@wrindexentrysplit\imki@wrindexentrysplit
\renewcommand\imki@wrindexentrysplit[3]{%
  \imki@orig@wrindexentrysplit{#1}{#2|hyperpage}{#3}%
}

\makeatletter
% --- Section titles: capture + one bookmark per logical section -------------
% \section is redefined so that:
%   1. the section title is captured into the plain-TeX register
%      \mytalksectitletoks (reprinted verbatim on the auto divider frame,
%      see the cmd/section/after hook below -- plain-TeX storage is
%      catcode-safe, unlike expanding the title in a hook), and
%   2. exactly ONE depth-1 PDF bookmark is created per \section call
%      (hyperref's own bookmark is suppressed via \Hy@writebookmark).
\DeclareCommandCopy{\OriginalSectionForBookmarks}{\section}
\newtoks\mytalksectitletoks

\RenewDocumentCommand{\section}{s o m}{%
  \begingroup
    \def\Hy@writebookmark##1##2##3##4##5{}%
    \global\mytalksectitletoks={#3}%
    \IfBooleanTF{#1}
      {\phantomsection\IfNoValueTF{#2}{\OriginalSectionForBookmarks*{#3}}{\OriginalSectionForBookmarks*[#2]{#3}}}%
      {\IfNoValueTF{#2}{\OriginalSectionForBookmarks{#3}}{\OriginalSectionForBookmarks[#2]{#3}}}%
  \endgroup
  \IfBooleanTF{#1}
    {\IfNoValueTF{#2}{\pdfbookmark[1]{#3}{section-star.\thepage.\number\inputlineno}}{\pdfbookmark[1]{#2}{section-star.\thepage.\number\inputlineno}}}%
    {\IfNoValueTF{#2}{\pdfbookmark[1]{Section \thesection: #3}{section.\thesection}}{\pdfbookmark[1]{Section \thesection: #2}{section.\thesection}}}%
}
\makeatother

% --- One PDF bookmark per logical frame (NOT per overlay) ---------------------
% ltx-talk writes one PDF page per overlay, so a naive frametitle hook would
% register a new bookmark for every overlay of the same frame. The
% \g_slideutils_frame_bookmark_int register remembers the last frame that
% got a bookmark, so subsequent overlays of the same frame are skipped.
% This is also what keeps BeamerPresenter's cumulative annotations working:
% one logical frame must stay one logical unit in the PDF outline.
\ExplSyntaxOn
\int_new:N \g_slideutils_frame_bookmark_int

\AddToHook{cmd/frametitle/after}{%
  \tl_if_empty:NF \g__talk_frame_title_tl
    {
      \ifnum\value{frame}=\g_slideutils_frame_bookmark_int\relax
      \else
        \int_gset:Nn \g_slideutils_frame_bookmark_int { \value{frame} }%
        \pdfbookmark[2]{\tl_use:N \g__talk_frame_title_tl}
          {frame.\int_use:N \g_slideutils_frame_bookmark_int}
      \fi
    }
}
\ExplSyntaxOff
\makeatother

% --- Printed index + theme for the colored boxes -------------------------------
% \makeindex[...]: two-column index titled "Index", listed in the TOC.
% The two EditInstance calls style the standard header/footer instances
% of talkthemetcolorbox (cyan background; footer order title/author/
% date/framenumber).
\makeindex[columns=2, title=Index, intoc]\EditInstance{header}{std}{
  background-color = cyan,
  color            = black,
  height           = 0.75cm
}
\EditInstance{footer}{std}{
  background-color = cyan,
  element-order = {title, author, date, framenumber},
}

% --- Auto section divider frame ---------------------------------------------------
% After each \section, insert a centered, large, structure-colored frame
% that prints the section title (from \mytalksectitletoks).
\AddToHook{cmd/section/after}{%
  \begin{frame}
    \centering
    \color{structure}
    \LARGE
    \the\mytalksectitletoks
  \end{frame}%
}

% --- Tagging sockets (tagged PDF structure) ---------------------------------------
% ltx-talk exposes tagging "sockets" where the PDF structure can attach
% section titles; register one and leave it unassigned (none) so the
% class's own tagging keeps control.
\ExplSyntaxOn
\AtBeginDocument{
  \NewTaggingSocketPlug { talk / sec / title } { none } { }
  \AssignTaggingSocketPlug { talk / sec / title } { none }
}
\ExplSyntaxOff

% --- Accessibility roles ----------------------------------------------------------------
% Map \frametitle to the H1 heading role in the tagged PDF, so screen
% readers treat slide titles as headings.
\tagpdfsetup{role / new-tag = frametitle / H1}

\makeatletter
% --- Page label = frame number -------------------------------------------------------
% Redefine what feeds the PDF page label to track the frame counter,
% not the physical page counter. Replace \theframenumber below with
% whatever internal macro/property ltx-talk actually uses to produce
% the "framenumber" footer element — confirm via ltx-talk-frame.dtx.
\renewcommand\thepage{\arabic{frame}}
\makeatother

\addbibresource{refs.bib}

% --- "Problem Time!" slides ------------------------------------------------------------------
% \ProblemSlide{problem}{solution} emits two frames: the problem, then
% its solution after a \pause, each numbered by the problem counter.
\newcounter{problemcounter}
\setcounter{problemcounter}{0}
\newcommand{\ProblemSlide}[2]{
  \stepcounter{problemcounter}
  \begin{frame}
    \frametitle{Problem Time!}
    Problem \theproblemcounter: #1
  \end{frame}

  \begin{frame}
    \frametitle{Problem Time!}
    Solution to Problem \theproblemcounter: \pause

    #2
  \end{frame}
}
% \highlight[<colour>]{<stuff>}
% from https://tex.stackexchange.com/questions/33401/a-version-of-colorbox-that-works-inside-math-environments
\newcommand{\highlight}[2][yellow]{\mathchoice%
  {\colorbox{#1}{$\displaystyle#2$}}%
  {\colorbox{#1}{$\textstyle#2$}}%
  {\colorbox{#1}{$\scriptstyle#2$}}%
  {\colorbox{#1}{$\scriptscriptstyle#2$}}}%

% --- Fonts ------------------------------------------------------------------------------
% Lato (text) + LeteSansMath (math): open, accessible typefaces.
\setmainfont{Lato}[
  Ligatures=TeX,
  Numbers={Lining,Monospaced}
]
\setsansfont{Lato}[
  Ligatures=TeX,
  Numbers={Lining,Monospaced}
]
\setmathfont{LeteSansMath.otf}\usepackage{emoji}

% --- Accessible colors ---------------------------------------------------------------------
% WCAG-contrast-safe variants of red/green for text on white:
% red!80!black and green!40!black both clear the 4.5:1 requirement.
% AltYellow: soft background for highlight boxes.
\colorlet{AccessibleRed}{red!80!black}
\colorlet{AccessibleGreen}{green!40!black}
\definecolor{AltYellow}{RGB}{255,245,170}

% --- Links ---------------------------------------------------------------------------------
% hyperref + link colors: every link uses AccessibleGreen (7.2:1 on
% white), so internal, citation, and URL links all pass WCAG 1.4.3.
\usepackage{hyperref}
\hypersetup{
  colorlinks=true,
  linkcolor=AccessibleGreen,
  urlcolor=AccessibleGreen,
  citecolor=AccessibleGreen,
}

% --- Shared course macros ---------------------------------------------------------------------
% Root macros.tex: math shorthands used across all course documents
% (\RR, \NN, \ZZ, \QQ, \grad, \curl, \tr, \der, \pder, ...).
\newcommand{\RR}{\mathbb{R}}
\newcommand{\NN}{\mathbb{N}}
\newcommand{\ZZ}{\mathbb{Z}}
\newcommand{\QQ}{\mathbb{Q}}
\newcommand{\vv}{\vec{v}}
\newcommand{\ww}{\vec{w}}
\newcommand{\xx}{\vec{x}}
\newcommand{\yy}{\vec{y}}
\newcommand{\zz}{\vec{z}}
\newcommand{\uu}{\vec{u}}
\newcommand{\basis}[1]{\vec{e_{#1}}}
\newcommand{\id}{\mathrm{id}}
\newcommand{\CC}{\mathbb{C}}
\renewcommand{\div}{\operatorname{div}}
\newcommand{\grad}{\operatorname{grad}}
\newcommand{\curl}{\operatorname{curl}}
\newcommand{\tr}{\operatorname{tr}}
\newcommand{\der}[2]{\dfrac{d #1}{d #2}}
\newcommand{\pder}[3][]{\dfrac{\partial^{#1} #2}{\partial #3}}


% % Title and author(s)
%%%%%%%%%%%%%%%%%%%%%%%%%%%%%%%%%%%%%%%%%%%%%%%%%%%%%%%
\title{Lecture 7 - Interpreting, Estimating, and Applying the 1st Derivative and the 2nd Derivative}
\author{Garrett Credi}
\date{09/29/26}
%%%%%%%%%%%%%%%%%%%%%%%%%%%%%%%%%%%%%%%%%%%%%%%%%%%%%%%
% %
% % The next command allows your in import encapsulated
% % postscript files, .epsf or .eps files, which
% % contain vector graphic image data.
% %
%%%%%%%%%%%%%%%%%%%%%%%%%%%%%%%%%%%%%%%%%%%%%%%%%%%%%%%
% %   The first thanks indicates your affiliation
% %
% %  Just the name here.
% %
% % Your mailing address goes at the end.
% %
% % \thanks is also how you indicate grant support
% %
%%%%%%%%%%%%%%%%%%%%%%%%%%%%%%%%%%%%%%%%%%%%%%%%%%%%%%%

%%%%%%%%%%%%%%%%%%%%%%%%%%%%%%%%%%%%%%%%%%%%%%%%%%%%%%%
% $ The next command ensures the margins are 1 inch on
% % all sides.
% %
%%%%%%%%%%%%%%%%%%%%%%%%%%%%%%%%%%%%%%%%%%%%%%%%%%%%%%%
\begin{document}
\begin{frame}
  \maketitle
\end{frame}

\begin{frame}
  \frametitle{Announcements}

  \textcolor{red}{Next week, during class time we will have our first exam! It will cover material from sections 1.1-1.7 and 3.2.3 of Active Calculus. \pause
  I.e. all the lectures up to and including this Thursday's lecture. \pause}

  Practice exams and the review sheet are present in the Canvas. \pause

  How to study? \pause
  \begin{itemize}
    \item Practice, practice, and more practice! The homework problems are your friends - use them to get practice solving the relevant problems. \pause
    \item Lecture Slides for review. I'll upload a compiled version of all the slides with a searchable index at the bottom so you can quickly go to 
      the slides defining relevant terms. \pause
    \item Use the practice exams as a \textit{cumulative metric} not as an additional source of problems. Take them after you complete a round of review,
      actually take them in an exam-like environment (no notes, 50 minutes, uncomfy chair, etc), and use your performance on them to guide your next round of review. \pause
    \item Start early! The later you start, the less time you'll have to ask me questions in office hours/via email! \pause
    \item Review session - I'll hold an Exam 1 review session next Monday (from 6-8?) where you can come and ask any and all questions you have and I'll answer em! 
  \end{itemize}
\end{frame}

\begin{frame}
  \frametitle{Quiz Time!}

  Please get into neat rows, and spread out. Only then will I hand out the quiz. As per usual, no cheating (hopefully should go without saying).
\end{frame}

\section{Interpreting and Applying the First Derivative}

\begin{frame}
  \frametitle{Derivative ReCap}    

  Last week we studied the \textit{derivative} of a function $f(x)$: \pause
  \begin{itemize}
    \item At a point $x = a$, we defined it as \pause
      \begin{equation}
        f'(a) = \lim\limits_{h \rightarrow 0} \dfrac{f(a + h) - f(a)}{h} = \lim\limits_{x \rightarrow a} \dfrac{f(x) - f(a)}{x - a}
      \end{equation} \pause
    \item We saw that it was also equal to: \pause
      \begin{enumerate}
        \item The slope of the best linear approximation to $f(x)$ at $x = a$ \pause
        \item The instantaneous velocity of a particle whose position was given by $f(x)$ at $t = a$ \pause
        \item The slope of the line tangent to the graph of $y = f(x)$ at $(a, f(a))$. \pause
      \end{enumerate}
    \item We also saw that we could construct a \textit{derivative function}: \pause
      \begin{equation}
        f'(x) = \lim_{h \rightarrow 0} \dfrac{f(x+h) - f(x)}{h}
      \end{equation} \pause
      and determine some information about $f(x)$ from $f'(x)$.
  \end{itemize}
\end{frame}

\begin{frame}
  \frametitle{Increasing-ness, Decreasing-ness, and Stationary-ness I}

  We can use the derivative to determine whether a function is increasing/decreasing/stationary: \pause

  \begin{block}{Increasing Function}\index{function!increasing}\index{increasing function}
    A function $f(x)$ is increasing around a point $x = a$
    if it is increasing in \textit{some} small neighborhood around $x = a$. \pause
  \end{block}

  Decreasingness is similarly defined. \pause

  \begin{alertblock}{Increasingness via the Derivative}\index{derivative!increasingness}\index{function!increasing!derivative test}
    For a function $f(x)$ which is differentiable at $x = a$, \pause
    increasingness at $a$ is equivalent to $f'(a) > 0$.

    Similarly, decreasingness is equivalent to $f'(a) < 0$.
  \end{alertblock}
\end{frame}

\begin{frame}
  \frametitle{Increasingness II}

  Why is that true? \pause

  Idea: Because $f(x)$ is differentiable at $x = a$, \pause
  it's behavior is well approximated by the best linear approximation
  $\mathscr{L}(x) = f(a) + f'(a)(x - a)$. \pause

  When is a linear function increasing/decreasing? \pause When the slope (which is $f'(a)$) is positive/negative! \pause

  In terms of instantaneous velocity: \pause positive velocity means increasing position, resp. negative velocity means decreasing position.
\end{frame}

\begin{frame}
  \frametitle{Graph Sketching I}

  We can use this idea to sketch graphs of derivative functions $f'(x)$
  given graphs of $f(x)$: \pause
  \vfill \pause

  And to sketch graphs of function $f(x)$ given graphs of their derivatives $f'(x)$:
  \vfill 
\end{frame}

\begin{frame}
  \frametitle{Graph Sketching II}

  Notice how the increasingness/decreasingness of the \textit{derivative} (e.g. the derivative of the derivative) informed us about the shape of the original function $f(x)$ \pause - we'll see later in lecture this is due to the properties of the \textit{second} derivative $f''(x)$.
\end{frame}

\section{Estimating The Derivative}

\begin{frame}
  \frametitle{Estimation I}
  Oftentime, in the real world, we won't have a literal function
  of which to take the derivative \pause we'll instead have some (possibly large) collection of sampled datapoints. \pause

  How can we get information about the derivative of the function modeling the measured quantity from these sparse datapoints? \pause

  The derivative is a limit so: \pause 
  \begin{equation}
    \text{Estimating the derivative} \quad \iff \quad \text{Estimating a limit!} 
  \end{equation} \pause
\end{frame}
\begin{frame}
  \frametitle{Estimation II}
  If we're trying to estimate: \pause
  \begin{equation}
    f'(a) = \lim\limits_{x \rightarrow a} \dfrac{f(x) - f(a)}{x - a} = \lim\limits_{h \rightarrow 0} \dfrac{f(a + h) - f(a)}{h}
  \end{equation} \pause

  We can compute the difference quotient quantities for values of $x \approx a$ or $h \approx 0$. \pause

  However, practically, there is a better choice of points about which to approximate: \pause

  \begin{block}{Central Difference Approximation}\index{central difference approximation}\index{derivative!central difference approximation}
    For a collection of sampled input and output pairs $(x_i, y_i)$, \pause
    estimating the derivative of $y$ with respect to $x$ about $x = a$, $\dfrac{dy}{dx} \mid_{x = a}$ \pause
    is best achieved by taking a pair of points $(x_0, y_0)$ and $(x_1, y_1)$ so that: \pause
    \begin{enumerate}
      \item $x_0 \approx a$ and $x_1 \approx a$ (approximation requirement) \pause
      \item $x_0$ and $x_1$ are ``symmetric" about $x = a$, \pause i.e. one of them is close and slightly less than $a$ while the other is close and slightly greater than $a$. \pause
    \end{enumerate}
  \end{block}
\end{frame}

\begin{frame}
  \frametitle{Central Difference Approximation Example} 

  E.g. \pause
  \vfill
\end{frame}

\section{Second and Higher Derivatives}

\begin{frame}
  \frametitle{Graph Sketching Redux}

  Recall that when we sketched graphs of $f(x)$ given the graph of $f'(x)$ \vfill \pause

  in addition to the positivity/negativity of $f'(x)$ giving us information about the increasingness/decreasingness of $f(x)$ \pause
  the increasingness/decreasingness of $f'(x)$ itself also gave us information \pause about the curvature of $f(x)$! \pause
\end{frame}

\begin{frame}
  \frametitle{The Second Derivative}
  The quantity that controls the increasingness/decreasingness of the \textit{first} derivative $g(x) := f'(x)$ \pause is \textit{it's} derivative: $g'(x) \pause = f''(x)$. \pause

  \begin{block}{Second Derivative}\index{derivative!second derivative}\index{second derivative}
    For a function $f(x)$ that is differentiable at $x = a$ \pause
    and so that $f'(x)$ is \textit{continuous} at $x = a$ \pause (which we say $f(x)$ is continuously differentiable at $x = a$) \pause,
    the second derivative of $f(x)$ at $x = a$ denoted: \pause
    \begin{equation}
      f''(x) \quad f^{(2)}(x) \quad \dfrac{d^2 f}{dx^2} \quad \dfrac{d^2}{dx^2} f(x)
    \end{equation} \pause

    is defined to be the derivative of the function $f'(x)$ at $x = a$. \pause

    The second derivative function is defined similarly.
  \end{block}
\end{frame}

\begin{frame}
  \frametitle{Quadratic Approximation}
  Just as the first derivative controls information about the best \textit{linear approximation}: \pause
  \begin{equation}
    f(x) \approx f(a) + f'(a)(x - a)
  \end{equation}\pause

  The second derivative controls information about the best \textit{quadratic approximation}: \pause
  \begin{equation}
    f(x) \approx f(a) + f'(a)(x - a) + \dfrac{f''(a)}{2} (x-a)^2
  \end{equation} \pause

  Similarly, just as the first derivative is equivalent to instantaneous \textit{velocity} \pause, the second derivative is equivalent to instantaneous \textit{acceleration}!
\end{frame}

\begin{frame}
  \frametitle{Concavity I}
  
  \begin{center}
    \begin{tabular}{c|c}
      $f(x) > 0$/$f(x) < 0$ & positivity/negativity \pause \\
      \hline
      $f'(x) > 0$/$f'(x) < 0$ & increasingness/decreasingness \pause \\
      \hline
      $f''(x) > 0$/$f''(x) < 0$ & ?
    \end{tabular}
  \end{center}
    \begin{block}{Concavity}\index{concavity}\index{function!concavity}
      A function $f(x)$ is \textit{concave up} at $x = a$ if,
      around $x = a$ \pause the function is \textit{greater} than it's 
      best linear approximation: \pause
      \begin{equation}
        f(x) > f(a) + f'(a)(x - a)
      \end{equation} \pause
      equivalently that it lies above it's tangent lines. \pause

      Similarly, it is \textit{concave down} at $x = a$ if: \pause
      \begin{equation}
        f(x) < f(a) + f'(a)(x - a)
      \end{equation} \pause
      i.e. that is lies below it's tangent lines.
    \end{block}
\end{frame}
\begin{frame}
  \frametitle{Concavity II}

  \begin{alertblock}{Concavity in terms of the Second Derivative}\index{concavity!in terms of the second derivative}
    A function $f(x)$ which is second differentiable at $x = a$ \pause
    is concave up at $x = a$ if and only if $f''(a) > 0$. \pause

    Similarly it is concave down if and only if $f''(a) < 0$. \pause
  \end{alertblock}

  E.g. 
  \vfill
\end{frame}

\begin{frame}
  \frametitle{Classifying Shapes of Functions via the First and Second Derivatives}

  We generically have four types of shapes a function can take
  \pause 
  given information about it's first and second derivatives: \pause

  \begin{columns}
    \begin{column}{0.49\textwidth}
      If $f'(x) > 0$ and $f''(x) > 0$: \pause
      \vfill
    \end{column} 

    \begin{column}{0.49\textwidth}
      If $f'(x) < 0$ and $f''(x) > 0$: \pause
      \vfill
    \end{column}
  \end{columns}
  \vfill

  \begin{columns}
    \begin{column}{0.49\textwidth}
      If $f'(x) > 0$ and $f''(x) < 0$: \pause
      \vfill
    \end{column}

    \begin{column}{0.49\textwidth}
      If $f'(x) < 0$ and $f''(x) < 0$: \pause
      \vfill
    \end{column}
  \end{columns}
  \vfill
\end{frame}

\begin{frame}
  \frametitle{Worksheet Time!}

  Let's get into groups and work on the worksheet!
\end{frame}

\begin{frame}
  \frametitle{End of Class}
  That's all folks!
\end{frame}
\end{document}
